\documentclass[11pt,a4paper]{article}

\usepackage[T1]{fontenc}
\usepackage{lmodern}
\usepackage{textcomp}
\usepackage{amsmath,amssymb,amsthm}
\usepackage{mathtools}
\usepackage{graphicx}
\usepackage{array}
\usepackage{geometry}
\usepackage{hyperref}
\usepackage{doi}
\geometry{margin=2.5cm}
\setlength{\emergencystretch}{2em}
\hypersetup{
  pdftitle={Independent Effective Subsystems under Non-Injective Admissibility: Exact Support and Preparation Criteria},
  pdfauthor={Jérôme Beau},
  pdfsubject={Support and preparation criteria for emergent effective subsystem composition},
  pdfkeywords={non-injective projection, effective subsystems, emergent composition, bicliques, rank one, steering}
}

\input{glyphtounicode}
\pdfgentounicode=1

\graphicspath{{out/}{../out/}}

\newtheorem{theorem}{Theorem}[section]
\newtheorem{lemma}[theorem]{Lemma}
\newtheorem{corollary}[theorem]{Corollary}
\newtheorem{proposition}[theorem]{Proposition}
\newtheorem{definition}[theorem]{Definition}
\newtheorem{remark}[theorem]{Remark}
\newtheorem{example}[theorem]{Example}

\newcommand{\Om}{\Omega}
\newcommand{\Aref}{\mathcal{A}_0}
\newcommand{\AC}[1]{\mathcal{A}_{#1}}
\newcommand{\OA}{O_A}
\newcommand{\OB}{O_B}
\newcommand{\PiA}{\Pi_A}
\newcommand{\PiB}{\Pi_B}
\newcommand{\PiAB}{\Pi_{AB}}
\newcommand{\SAB}{\mathcal{S}_{AB}}
\newcommand{\SAloc}{\mathcal{S}_A^{\mathrm{loc}}}
\newcommand{\SBloc}{\mathcal{S}_B^{\mathrm{loc}}}
\newcommand{\rk}{\operatorname{rank}}

\title{Independent Effective Subsystems under Non-Injective Admissibility:\\
Exact Support and Preparation Criteria}

\author{J\'{e}r\^{o}me Beau\\
\small Independent Researcher}

\date{2026\\[2pt]\small Version 1.0.0}

\begin{document}

  \maketitle

  \begin{abstract}
    A recurring assumption in emergent-physics programmes is that once a fundamental description is given,
    the decomposition of the world into independently preparable subsystems comes for free.
    We show that on a substrate not supplied with a subsystem factorisation this assumption splits into two
    separate, checkable layers, and we characterise both exactly in the finite case.
    The arena is a finite configuration set $\Om$ carrying a reference admissible subset, two effective
    projections $\PiA \colon \Om \to \OA$ and $\PiB \colon \Om \to \OB$ that are not assumed injective, and
    cylindrical admissibility conditions combined by intersection; no factorisation
    $\Om \simeq \Om_A \times \Om_B$ is assumed at any point.
    At the possibilistic layer, we prove that the pairs of local outcome sets on which joint interrogation
    is margin-stable and jointly rectangular --- the \emph{independence blocks} --- are exactly the nonempty
    rectangles (bicliques) contained in the projected support relation.
    At the preparation layer, we prove a benchmark and a criterion.
    The benchmark: on any block, if all probability measures on the underlying support are admissible
    preparations, every joint outcome law, in particular every product law, lifts to the substrate; fundamental
    the absence of a supplied substrate factorisation alone therefore obstructs nothing.
    The criterion: for the family of normalised cylindrical conditionings of a strictly positive reference
    weight matrix $W$, closure under products of independently locally conditioned states holds if and only
    if $W$ has rank one.
    A minimal countermodel exhibits a complete rectangular support whose induced preparations are
    nevertheless correlated: support rectangularity does not imply preparation independence.
    A type discipline runs through both layers: states of one side obtained by conditioning on the other are
    steered conditional states, not local preparations, and admitting them manufactures spurious witnesses.
    The individual mathematical ingredients are elementary or standard; the contribution is their typed
    assembly into a necessary support-and-preparation contract for emergent effective subsystems over an
    admissibility structure without a supplied subsystem factorisation.
    The hypotheses abstract the composition problem motivated by the non-injective foundations programme,
    but their realisation from the programme's axioms is not derived here.
    \emph{Interpretive reading (structural, not a result):} the absence of a fundamental subsystem
    factorisation does not forbid independent effective subsystems; it relocates their independence from an
    inherited property of the substrate to an emergent property of the projected support and of the
    preparation weights.
  \end{abstract}

  \paragraph{Keywords.}
    non-injective projection; effective subsystems; emergent composition; possibilistic independence;
    bicliques; rank-one criterion; steering; local tomography.

    \clearpage
    \tableofcontents

    \clearpage
  \section{Introduction}
  \label{sec:intro}

  \subsection{The composition gap}
    \label{ssec:gap}

    The non-injective foundations programme~\cite{Beau2026m} describes observable physics as
    the image of a many-to-one projection from a relational substrate.
    Non-injectivity is structurally forced for genuine emergence~\cite{Beau2026l}, and its observable
    signatures include Bell-type correlations obtained without any subsystem decomposition of the underlying
    configuration space~\cite{Beau2026b}.
    The substrate of that argument is not supplied with a subsystem factorisation: no split
    $\Om \simeq \Om_A \times \Om_B$ is imposed, and the measurement regions exist only at the level of
    observable descriptions.

    This creates a gap on the other side of the ledger.
    Laboratory physics is organised around \emph{independent subsystems}: systems that can be addressed
    separately, prepared independently, and recombined, with joint statistics generated from local ones.
    If the substrate is not supplied with a subsystem factorisation, none of this is given.
    The present paper asks the resulting question in its sharpest finite form:

    \begin{quote}
      On a substrate without a supplied subsystem factorisation, equipped with effective projections that
      are not assumed injective, when do independently addressable and preparable effective subsystems exist
      --- and what exactly must be added, at the support and preparation levels, for independent composition
      to be possible?
    \end{quote}

  \subsection{A missing transversal contract}
    \label{ssec:k4}

    The corpus already consumes an answer it does not yet possess.
    The universal singlet construction builds bipartite proto-states inside a tensor product
    $\chi_{2j+1} \otimes \chi_{2j+1}$ of representation spaces~\cite{Beau2026q3}, and the entanglement of
    conjugate Weil pairs is defined on a bipartite space $W_c(n) \otimes W_{q-c}(n)$ introduced by an explicit
    closure hypothesis~\cite{Beau2026entnote}.
    In both cases the composite structure is supplied, not derived; these works are consumers of a
    composition structure, and are cited here as such, not as evidence that the structure exists.

    We name the missing piece as a transversal contract:
    \begin{equation}
      \label{eq:k4}
      \mathrm{K4}\colon\quad
      \mathsf{Admissibility}
      \;\longrightarrow\;
      \mathsf{Composable\ effective\ subsystems}.
    \end{equation}
    The arrow is the open bridge; its source carries no supplied subsystem factorisation.
    This paper closes two preconditions of~\eqref{eq:k4} --- the support layer and the preparation layer ---
    and leaves open its states-and-effects layer: the construction of effect spaces, their bilinear pairing
    with states, local tomography, and any quantum tensor product.

  \subsection{Contributions}
    \label{ssec:contributions}

    All results are finite and exact, and are validated by exhaustive enumeration scripts distributed with
    the paper (Appendix~\ref{app:repro}).

    \begin{enumerate}
      \item \textbf{Block theorem (Theorem~\ref{thm:blocks}).}
      Under cylindrical conditions with conjunction by intersection, the local outcome-set pairs passing the
      two composability tests --- margin stability and joint rectangularity --- are exactly the nonempty
      rectangles $S_A \times S_B \subseteq R$ of the projected support relation $R$; maximal blocks are the
      maximal bicliques of $R$.
      \item \textbf{Lifting benchmark (Theorem~\ref{thm:lifting}).}
      On a block, if every probability measure on the underlying support is an admissible preparation, the
      push-forward to joint outcome laws is surjective; every product of local states is jointly preparable.
      The absence of a supplied substrate factorisation alone therefore obstructs nothing at this layer.
      \item \textbf{Rank-one criterion (Theorem~\ref{thm:rankone}).}
      For the family of normalised cylindrical conditionings of a strictly positive weight matrix $W$ on a
      finite block, closure under products of independently locally conditioned states is equivalent to
      $\rk W = 1$.
      \item \textbf{Separation countermodel (Example~\ref{ex:n3}).}
      A $2 \times 2$ block with weights $N_3 = \begin{psmallmatrix} 1 & 1 \\ 1 & 2 \end{psmallmatrix}$ has
      complete rectangular support and a non-product-closed conditioning family: support rectangularity does
      not imply preparation independence.
      \item \textbf{Steering discipline (Definition~\ref{def:local-families}).}
      Local state families must be generated by local conditionings only.
      Margins of joint conditionings are steered conditional states; admitting them as local preparations
      changes the witness count of the countermodel and manufactures spurious failures.
    \end{enumerate}

    The individual mathematical ingredients are elementary or standard; the contribution is their typed
    assembly into a necessary support-and-preparation contract for emergent effective subsystems over an
    admissibility structure without a supplied subsystem factorisation.

  \subsection{What this paper does not claim}
    \label{ssec:nonclaims}

    This paper derives no quantum tensor product, Born rule, quantum field, or local-tomography property.
    The arena of Section~\ref{sec:arena} abstracts the composition problem motivated by the foundations
    programme, but its realisation from the programme's axioms is not derived here; the hypotheses are a
    contract, stated as such.
    The reference measure of Section~\ref{sec:preparation}, its conditioning rule, and the physical
    availability of convex mixing are construction choices.
    The rank-one criterion is a criterion within the cylindrical-conditioning class; other admissible
    preparation families require their own criteria.
    All statements are finite; the strictly positive case carries the converse of
    Theorem~\ref{thm:rankone}.

  \section{The Arena: Substrate, Projections, and Conditional Preparations}
  \label{sec:arena}

  \subsection{Substrate and effective projections}
    \label{ssec:substrate}

    Let $\Om$ be a finite set of configurations and $\Aref \subseteq \Om$ a reference admissible subset.
    Every device, setting, and boundary condition relevant to the situations described below is part of the
    same $\Om$; combining two conditions never compares different spaces.

    Two maps
    \begin{equation}
      \PiA \colon \Om \to \OA,
      \qquad
      \PiB \colon \Om \to \OB
    \end{equation}
    are the effective projections of the two prospective subsystems, and
    $\PiAB = (\PiA, \PiB) \colon \Om \to \OA \times \OB$ is their joint application.
    The labels $A$ and $B$ name independently addressable effective subsystems, not spatial regions; the
    construction precedes any geometry, and spatial regions would arise, if at all, as a later special case.
    No physically distinguished factorisation $\Om \simeq \Om_A \times \Om_B$ is assumed, and none will be
    used.
    The product appears only on the observable side, as the codomain of a pair of maps out of one undivided
    substrate.
    Accordingly, ``without a supplied factorisation'' means that no such identification or subsystem maps
    belong to the physical data; it does not claim that the finite set $\Om$ admits no accidental bijection
    with some Cartesian product.

  \subsection{Cylindrical conditions and the meet}
    \label{ssec:conditions}

    \begin{definition}[Cylindrical condition]
      \label{def:cylindrical}
      A local possibilistic condition on side $A$ is a subset of allowed outcomes $S_A \subseteq \OA$,
      realised as
      \begin{equation}
        \AC{C_A} = \Aref \cap \PiA^{-1}(S_A),
      \end{equation}
      and symmetrically for side $B$.
    \end{definition}

    Local conditions of $A$ bear only on $A$, and those of $B$ only on $B$.
    The conjunction of two conditions is their intersection:
    \begin{equation}
      \label{eq:meet}
      \AC{C_A \wedge C_B} = \AC{C_A} \cap \AC{C_B}.
    \end{equation}
    The meet~\eqref{eq:meet} is commutative, associative, idempotent, and monotone (a strengthened condition
    never enlarges the admissible set); it is partial (an empty intersection means incompatibility of the
    two conditions, not an error); and it is extensional (two conditions selecting the same subset of $\Om$
    are equivalent at this level).
    If combining two devices changes admissibility other than by intersection, the meet is not widened:
    the interaction is modelled by an explicit additional context or by enlarging $\Om$, so that
    contextuality is never hidden inside the definition of~$\wedge$.

  \subsection{The two composability tests}
    \label{ssec:tests}

    Fix condition sets $S_A \subseteq \OA$, $S_B \subseteq \OB$ and write
    $\mathcal{M} = \AC{C_A} \cap \AC{C_B}$ for the meet support.
    Independent composability of the two selections requires two properties, which we deliberately separate
    because they fail differently.

    \begin{definition}[Test 1: margin stability]
      \label{def:test1}
      $\PiA(\mathcal{M}) = \PiA(\AC{C_A})$ and $\PiB(\mathcal{M}) = \PiB(\AC{C_B})$: the condition imposed
      at $B$ deletes no locally possible outcome at $A$, and conversely.
    \end{definition}

    \begin{definition}[Test 2: joint rectangularity]
      \label{def:test2}
      $\PiAB(\mathcal{M}) = \PiA(\mathcal{M}) \times \PiB(\mathcal{M})$: every pair of outcomes that remains
      locally possible is jointly possible.
    \end{definition}

    Together the two tests are equivalent to the single equality
    $\PiAB(\AC{C_A \wedge C_B}) = \PiA(\AC{C_A}) \times \PiB(\AC{C_B})$; separated, they diagnose why it
    fails: cross restriction of margins, correlation of the joint support, or both.

  \subsection{The bipartite reduction}
    \label{ssec:reduction}

    At the possibilistic level the whole structure is the relation
    \begin{equation}
      R = \PiAB(\Aref) \subseteq \OA \times \OB,
    \end{equation}
    with coordinate projections $\pi_A(R)$ and $\pi_B(R)$.
    Extensional representatives of conditions are subsets $S_A \subseteq \pi_A(R)$ and
    $S_B \subseteq \pi_B(R)$, the selected support is $R \cap (S_A \times S_B)$, and Test~2 asks whether the
    selected bipartite subgraph is a complete biclique.
    This reduction permits exhaustive verification on small examples with no probabilities, amplitudes, or
    tensors (Appendix~\ref{app:repro}).

  \section{The Possibilistic Layer: Independence Blocks are Bicliques}
  \label{sec:possibilistic}

  \begin{definition}[Independence block]
    \label{def:block}
    A pair $(S_A, S_B)$ of extensional condition sets is an \emph{independence block} of $R$ when Tests~1
    and~2 both hold with nonempty selected support.
  \end{definition}

  \begin{theorem}[Blocks are bicliques]
    \label{thm:blocks}
    $(S_A, S_B)$ is an independence block of $R$ if and only if
    $\emptyset \neq S_A \times S_B \subseteq R$.
  \end{theorem}

  \begin{proof}
    If $S_A \times S_B \subseteq R$, the selected support is
    $R \cap (S_A \times S_B) = S_A \times S_B$, whose margins are $S_A$ and $S_B$; both tests hold.
    Conversely, write $\mathrm{sel} = R \cap (S_A \times S_B)$.
    Test~2 gives $\mathrm{sel} = \pi_A(\mathrm{sel}) \times \pi_B(\mathrm{sel})$, and Test~1 identifies the
    margins as $\pi_A(\mathrm{sel}) = S_A$ and $\pi_B(\mathrm{sel}) = S_B$; hence
    $S_A \times S_B = \mathrm{sel} \subseteq R$.
  \end{proof}

  \begin{remark}[Order structure and formal concepts]
    \label{rem:fca}
    The nonempty independence blocks form the poset of nonempty rectangles contained in $R$, ordered by
    componentwise inclusion.
    Its maximal elements are the nonempty maximal bicliques of $R$, which correspond to the nondegenerate
    formal concepts of the context $(\OA, \OB, R)$ in the sense of formal concept
    analysis~\cite{GanterWille1999}.
    The formal-concept lattice, however, is built on extent inclusion together with Galois closure; it is
    not a lattice structure on all independence blocks under componentwise inclusion.
    In the L-shaped relation of Table~\ref{tab:relations}, the two maximal blocks have no join: the
    rectangle they generate contains the missing corner.
  \end{remark}

  \begin{remark}[Database reading]
    \label{rem:database}
    Cylindrical selection is the relational-algebra selection operator, and full rectangularity of a
    relation is the trivial binary join dependency: $R$ decomposes losslessly as the product of its
    projections~\cite{Fagin1977}.
    Theorem~\ref{thm:blocks} is, in that vocabulary, a classification of the sub-selections of a
    non-decomposable relation that are locally decomposable.
  \end{remark}

  \subsection{Five reference relations}
    \label{ssec:relations}

    Table~\ref{tab:relations} and Figure~\ref{fig:relations} report the exhaustive classification of five
    reference relations; the enumeration is reproduced by the committed script
    \texttt{code/possibilistic\_blocks.py}.

    \begin{table}[ht]
      \centering
      \footnotesize
      \setlength{\tabcolsep}{4.5pt}
      \begin{tabular}{l c c c c c}
        \hline
        Relation & pairs $(S_A, S_B)$ & Test~1 & Test~2 & blocks (nonempty) & maximal blocks \\
        \hline
        complete $2{\times}2$ & 16 & 10 & 16 & 9 & the full rectangle \\
        diagonal $2{\times}2$ & 16 & 4 & 15 & 2 & $\{0\}{\times}\{0\}$, $\{1\}{\times}\{1\}$ \\
        L-shape (three edges) & 16 & 7 & 15 & 5 & $\{0\}{\times}\{0,1\}$, $\{0,1\}{\times}\{0\}$ \\
        two disjoint $2{\times}2$ rectangles & 256 & 100 & 175 & 18 & the two rectangles \\
        6-cycle on $3{\times}3$ & 64 & 29 & 48 & 12 & six stars $K_{1,2}$ \\
        \hline
      \end{tabular}
      \caption{Exhaustive two-test classification over extensional condition pairs.
      Pair counts include the degenerate empty pair; block counts exclude it.}
      \label{tab:relations}
    \end{table}

    \begin{figure}[ht]
      \centering
      \includegraphics[width=\textwidth]{fig_relations}
      \caption{The five reference relations (grey edges) with their maximal independence blocks highlighted
      (coloured bicliques).
      Generated by \texttt{code/make\_figures.py}.}
      \label{fig:relations}
    \end{figure}

    The failure modes separate as follows.
    On the diagonal relation the dominant failure is cross restriction of margins (12 of 16 pairs fail
    Test~1); exactly one pair, the full selection, is margin-stable yet fails rectangularity.
    Only trivial one-outcome subsystems compose on a perfectly correlated support; this is the context-wise
    diagonal (or anti-diagonal) support pattern appearing in the Bell toy model of~\cite{Beau2026b}, not a
    general characterisation of Bell scenarios.
    The L-shape is the pure specimen of the second failure mode at full selection: margins stable, support
    correlated by the missing corner.
    The two disjoint rectangles produce a sector-like possibilistic decomposition: every block lives inside
    one rectangle, every cross-sector selection fails by margin cutting or incompatibility.
    The decomposition is analogous to a superselection pattern but is not an algebraic or operational
    superselection rule, which would concern states, observables, or algebras rather than a relational
    support.
    The 6-cycle contains no $K_{2,2}$, so every block is a star: composition is possible only when one side
    is pinned to a single outcome.

  \section{The Preparation Layer: Weights Decide What Support Cannot}
  \label{sec:preparation}

    The possibilistic layer controls which outcome pairs are jointly possible.
    It forgets two things: the multiplicities of the fibres $f^{-1}(a,b)$ of the restricted joint
    projection, and the weights that preparations place on them.
    This section shows that these forgotten data carry the whole question of preparation independence.

  \subsection{The simplex-completion benchmark}
    \label{ssec:benchmark}

    Fix an independence block $S_A \times S_B \subseteq R$ and write $X = \AC{C_A \wedge C_B}$ for its meet
    support and $f = \PiAB|_X \colon X \to S_A \times S_B$ for the restricted joint projection.
    On a block, $f$ is surjective: that is exactly what the biclique provides.

    \begin{theorem}[Lifting under simplex completion]
      \label{thm:lifting}
      If all probability measures on $X$ are admissible preparations, the push-forward
      $f_* \colon \Delta(X) \to \Delta(S_A \times S_B)$ is surjective.
      Explicitly, choosing one preimage $x(a,b) \in f^{-1}(a,b)$ for each pair, any joint law $\nu$ lifts as
      $\sum_{(a,b)} \nu(a,b)\, \delta_{x(a,b)}$, affinely in $\nu$ for a fixed section.
      In particular every product $\mu_A \otimes \mu_B$ with $\mu_A \in \Delta(S_A)$,
      $\mu_B \in \Delta(S_B)$ is realised by an admissible preparation.
    \end{theorem}

    \begin{proof}
      Immediate from surjectivity of $f$ on the block: the displayed lift is a probability measure on $X$
      and pushes forward to $\nu$.
    \end{proof}

    The symbol $\otimes$ above is the product of measures on a product of outcome sets, which exists
    set-theoretically on any block; it is not a tensor product of state spaces, and the composite state
    space is precisely what remains unconstructed.

    \begin{corollary}[Absence of a supplied factorisation alone obstructs nothing]
      \label{cor:noobstruction}
      Once a biclique is available, the absence of a supplied subsystem factorisation on $\Om$ does not by
      itself obstruct product preparations.
      Any obstruction comes from the restriction of the admissible preparation family.
    \end{corollary}

  \subsection{Preparation families and the steering discipline}
    \label{ssec:families}

    A preparation family over the block is a set $\mathcal{M}_{AB} \subseteq \Delta(X)$, with effective
    image $\SAB = (\PiAB)_* \mathcal{M}_{AB}$.
    Two control families calibrate the scale.
    The deterministic family (Dirac measures only) realises every deterministic product
    $\delta_a \otimes \delta_b$ on a block while offering no mixtures: preparability of outcomes and
    preparability of mixtures are different properties.
    The fibre-uniform family $L(\nu)(x) = \nu(a,b)/|f^{-1}(a,b)|$ for $x \in f^{-1}(a,b)$ is convex and
    invariant under internal fibre permutations, yet satisfies $f_* L(\nu) = \nu$ identically: its effective
    image is the full simplex.
    Uniformity within fibres constrains microscopic lifts and constrains effective preparation weights not
    at all.

    The genuinely restrictive family used below conditions a reference measure.
    Fix a positive measure $\lambda$ on $X$ and define, for nonempty $U \subseteq S_A$,
    $V \subseteq S_B$ with fibre-weight matrix $w_{ab} = \lambda(f^{-1}(a,b))$,
    \begin{equation}
      \label{eq:conditioning}
      p^{U,V}_{ab}
      \;=\;
      \frac{w_{ab}\, \mathbf{1}_U(a)\, \mathbf{1}_V(b)}
           {\sum_{a' \in U,\, b' \in V} w_{a'b'}}.
    \end{equation}
    Neither $\lambda$ nor the conditioning rule is derived from admissibility; both are construction
    choices, stated as such.

    \begin{definition}[Local state families]
      \label{def:local-families}
      The joint family is $\SAB = \{ p^{U,V} : \emptyset \neq U \subseteq S_A,\ \emptyset \neq V \subseteq S_B \}$.
      The \emph{local} state families use local conditions only:
      \begin{equation}
        \SAloc = \{ \pi_A\, p^{U, S_B} : \emptyset \neq U \subseteq S_A \},
        \qquad
        \SBloc = \{ \pi_B\, p^{S_A, V} : \emptyset \neq V \subseteq S_B \}.
      \end{equation}
    \end{definition}

    The restriction in Definition~\ref{def:local-families} is a type discipline, not a convenience.
    The margin $\pi_A\, p^{U,V}$ for $V \subsetneq S_B$ is a state of $A$ obtained by conditioning on $B$:
    a steered conditional state in the sense made operational for quantum
    assemblages~\cite{Schrodinger1935,Schrodinger1936,WisemanJonesDoherty2007}, transposed here to the
    classical-measure level.
    Steered states are not local preparations, and admitting them into $\SAloc$ changes theorems: on the
    countermodel of Example~\ref{ex:n3} it inflates the local family from three states per side to five and
    the failing pairs from five to seventeen, all the extra witnesses being illegitimate.

  \subsection{The rank-one criterion}
    \label{ssec:rankone}

    \begin{theorem}[Product closure for cylindrical conditionings]
      \label{thm:rankone}
      Let $W = (w_{ab})$ be strictly positive on the finite block $S_A \times S_B$.
      The family $\SAB$ is closed under products of independently locally conditioned states,
      \begin{equation}
        \forall\, \mu_A \in \SAloc,\ \forall\, \mu_B \in \SBloc \colon
        \quad
        \mu_A \otimes \mu_B \in \SAB,
      \end{equation}
      if and only if $\rk W = 1$, i.e.\ all $2 \times 2$ minors vanish:
      $w_{ab} w_{a'b'} = w_{ab'} w_{a'b}$.
    \end{theorem}

    \begin{proof}
      \emph{Direct direction.}
      If $w_{ab} = \alpha_a \beta_b$, then $p^{U,V} = \alpha^U \otimes \beta^V$, where $\alpha^U$ and
      $\beta^V$ are the normalised restrictions of $\alpha$ and $\beta$ to $U$ and $V$.
      Every locally conditioned state is such a restriction, and the product of two restrictions is the
      conditioning on the corresponding cylinder; every product therefore lies in $\SAB$.

      \emph{Converse.}
      The unconditioned margins $r = \pi_A\, p^{S_A, S_B}$ and $c = \pi_B\, p^{S_A, S_B}$ belong to
      $\SAloc$ and $\SBloc$.
      Closure demands $r \otimes c = p^{U,V}$ for some $U, V$.
      Since $W$ is strictly positive, $r \otimes c$ has full support $S_A \times S_B$, while the support of
      $p^{U,V}$ is $U \times V$; hence $U = S_A$ and $V = S_B$, so $p^{S_A,S_B} = r \otimes c$.
      The normalised law is then a product of its margins, which is equivalent to $\rk W = 1$.
    \end{proof}

    \begin{remark}
      \label{rem:nocoincidence}
      There is no accidental coincidence with another conditioning in the converse: full support excludes it
      immediately.
      Strict positivity is exactly the block condition for counting-type measures ($n_{ab} \geq 1$ on a
      biclique); for a general $\lambda$ it is an added hypothesis.
      The single-law direction of the criterion --- a law is a product of its margins iff its matrix has rank
      one --- is classical, and rank-one structure on a prescribed support is the quasi-independence model of
      contingency-table analysis~\cite{Goodman1968}.
      Theorem~\ref{thm:rankone} differs in statement type: it is a closure property of the entire
      conditioning family, quantified over locally conditioned states, with the full-support converse; the
      statement is elementary in substance and, to the author's knowledge, new as formulated.
    \end{remark}

  \subsection{The separation countermodel}
    \label{ssec:countermodel}

    \begin{example}[$N_3$: rectangular support, correlated preparations]
      \label{ex:n3}
      On the complete $2 \times 2$ biclique take counting weights
      \begin{equation}
        N_1 = \begin{pmatrix} 1 & 1 \\ 1 & 1 \end{pmatrix},
        \qquad
        N_2 = \begin{pmatrix} 1 & 2 \\ 2 & 4 \end{pmatrix},
        \qquad
        N_3 = \begin{pmatrix} 1 & 1 \\ 1 & 2 \end{pmatrix}.
      \end{equation}
      $N_1$ is uniform and factorising; $N_2$ is non-uniform and factorising (rank one); both families are
      product-closed, with three local states per side and nine conditioned laws.
      $N_3$ has minor $1 \neq 0$: its family fails closure on exactly five of the nine local pairs --- the
      four products of a deterministic state with the opposite unconditioned margin, and the product of the
      two unconditioned margins
      \begin{equation}
        \mu_A = \mu_B = \left( \tfrac{2}{5}, \tfrac{3}{5} \right),
        \qquad
        \mu_A \otimes \mu_B
        = \begin{pmatrix} 4/25 & 6/25 \\ 6/25 & 9/25 \end{pmatrix}.
      \end{equation}
      The product has full support; the only conditioned law with full support is the unconditioned law
      $\tfrac{1}{5} N_3$, and the two differ.
    \end{example}

    The two layers are therefore provably distinct:
    \begin{quote}
      the biclique controls the possibilistic zeros; the rank of $W$ controls the correlations of the
      chosen measure.
      Support rectangularity does not imply preparation independence.
    \end{quote}
    Non-uniformity is irrelevant ($N_2$); only rank matters.

  \section{Position within the Neighbouring Literature}
  \label{sec:position}

    The raw materials of this paper live in six developed literatures; we state for each what is shared,
    what differs, and why the present assembly is not subsumed.

    \paragraph{Formal concept analysis.}
    Maximal bicliques of a binary relation are the formal concepts of its context, organised by the Galois
    connection into the concept lattice~\cite{GanterWille1999}.
    Theorem~\ref{thm:blocks} borrows the objects, not the question: FCA studies the closure structure of a
    relation, while the two-test contract asks which cylinder selections compose, and the block poset under
    componentwise inclusion is not the concept lattice (Remark~\ref{rem:fca}).

    \paragraph{Database dependency theory.}
    Rectangularity is lossless join decomposability, and multivalued dependencies characterise
    it~\cite{Fagin1977}.
    The present use is local: a classification of the decomposable sub-selections of a globally
    non-decomposable relation, read as a composition arena rather than as a schema-design constraint.

    \paragraph{Contingency tables.}
    Rank-one is the classical independence criterion, and rank-one on a prescribed support is
    quasi-independence~\cite{Goodman1968}.
    The statistical literature fits and tests these as models of data; Theorem~\ref{thm:rankone} instead
    characterises closure of a preparation family, a statement about which mixtures are jointly realisable,
    with the steering discipline deciding which margins count as local.

    \paragraph{Generalized probabilistic theories.}
    Local tomography is a standard composition axiom in operational reconstructions and is among the
    conditions that distinguish complex quantum theory from real and quaternionic variants; generalised
    probabilistic theories take state and effect spaces, together with admissible composition rules, as
    operational input~\cite{Hardy2001,Barrett2007,ChiribellaDArianoPerinotti2011}.
    Recent work also shows that single-system quantum structure and Bell-like correlations do not uniquely
    determine the composition rule~\cite{ErbaPerinotti2024}.
    The present paper works one level below: it asks when independently addressable subsystems and product
    preparations exist at all, from projections on a substrate without a supplied factorisation, and treats
    the operational composite as the object to be earned rather than assumed.
    The corpus position that the GPT duality is a structure whose emergence is in question rather than a
    starting point~\cite{Beau2026l} is the mandate for that order of construction.

    \paragraph{Steering.}
    The distinction between states remotely selected by conditioning and states locally prepared goes back
    to Schr\"{o}dinger~\cite{Schrodinger1935,Schrodinger1936} and was made operational
    in~\cite{WisemanJonesDoherty2007}.
    Definition~\ref{def:local-families} transposes the cut to classical measures over a substrate, where it
    is already load-bearing: it changes the witness set of Example~\ref{ex:n3}.

    \paragraph{Contextuality and induced subsystem structure.}
    Possibilistic empirical models over measurement covers, and the sheaf-theoretic characterisation of
    contextuality as obstruction to global sections, share the raw object of
    Section~\ref{sec:possibilistic} --- a support relation --- but ask a different question of
    it~\cite{AbramskyBrandenburger2011}.
    Inside quantum theory, tensor product structures are induced by distinguished observable algebras
    rather than given~\cite{Zanardi2001,ZanardiLidarLloyd2004}; the present construction is a pre-Hilbert
    analogue in which even the possibilistic and convex prerequisites of a composite are conditional.
    The 6-cycle motif of Table~\ref{tab:relations} is the standard shape of possibilistic contextuality
    arguments; no link is asserted beyond the shared combinatorics.

  \section{Discussion}
  \label{sec:discussion}

  \subsection{Consequences for the programme}
    \label{ssec:consequences}

    The results are logically independent of the programme's carrier-selection and continuum chokepoints:
    no Heisenberg carrier, no Weil representation, no complex structure, and no geometry is consumed.
    They are compatible with any conditional realisation of those structures.
    What they contradict is a sufficiency reading: selecting a carrier does not by itself provide two
    independently preparable subsystems, a composition rule, a product-closed preparation family, an effect
    space, or local tomography; and a future spatial continuum would not by itself make disjoint regions
    independent subsystems, since the block condition and a preparation-closure condition would still have
    to hold.
    In the vocabulary of~\eqref{eq:k4}: the support and preparation preconditions of K4 are now
    characterised, its states-and-effects layer is untouched, and the corpus constructions that consume a
    composite~\cite{Beau2026q3,Beau2026entnote} retain their results conditionally on the supplied
    structure.

  \subsection{Interpretive outlook}
    \label{ssec:outlook}

    The following reading is interpretation --- a structural way of situating the results --- and not itself
    a result.
    The absence of a fundamental subsystem factorisation does not forbid independent effective subsystems.
    It relocates their independence: from an inherited property of the substrate to an emergent property of
    the projected support (which pairs of outcomes are jointly possible) and of the preparation weights
    (which mixtures are jointly realisable).
    What a laboratory calls ``two systems'' is, in this reading, a regime of the theory where both layers
    cooperate --- a biclique carrying a rank-one weight structure --- and the existence of such regions is a
    constraint on any future dictionary between the substrate description and operational physics, not a
    free assumption.

  \subsection{Limits and future work}
    \label{ssec:limits}

    The limits are those stated in Section~\ref{ssec:nonclaims}.
    Natural extensions, none of which affects the finite proofs above, include: weight matrices with zeros
    inside the block; approximate independence, with quantitative bounds in total variation, mutual
    information, or normalised minors when $\rk W \approx 1$; the measurable, non-finite case; the
    states-and-effects layer over a rank-one family, where local tomography first becomes statable for this
    construction; and the possible use of composed-probability coherence to constrain quantitative weight
    laws in the foundations programme.

  \section*{Acknowledgements}
    Portions of the analysis and editorial preparation benefited from iterative interactions with large
    language models used as analytical assistants.
    All claims and final formulations remain the sole responsibility of the author.

  \appendix

  \section{Reproduction}
  \label{app:repro}

    All numerical claims are reproduced by two deterministic scripts using exact rational arithmetic and no
    dependencies beyond the Python standard library:
    \begin{itemize}
      \item \texttt{code/possibilistic\_blocks.py} --- exhaustive two-test classification of the five
      relations of Table~\ref{tab:relations};
      \item \texttt{code/convex\_preparation\_tests.py} --- families and matrices of
      Section~\ref{sec:preparation}, with the local state families of Definition~\ref{def:local-families}.
    \end{itemize}
    Figure~\ref{fig:relations} is generated by the script \texttt{code/make\_figures.py}, under the pinned
    environment recorded in \texttt{code/requirements.txt}.
    From a fresh clone:
    \begin{verbatim}
python3 code/possibilistic_blocks.py
python3 code/convex_preparation_tests.py
bash build.sh
    \end{verbatim}
    The scripts print the full inventories; \texttt{build.sh} regenerates the figure, verifies it against
    \texttt{ARTIFACT\_SHA256SUMS}, and compiles the paper.

  \bibliographystyle{plain}
  \bibliography{cosmochrony-bibliography,external-refs}

\end{document}
